Small private medical chambers in low-resource settings still order patients by paper
tokens or verbal agreement, and they measure basic vitals inside the consultation room,
where clinician time is scarcest. We present a stand-alone embedded system that moves both
tasks onto a single low-cost microcontroller. Interaction is split across two separate
stations. A patient-side self-service booth takes identifying details and acquires vital
signs, while a doctor-side console shows the current patient's record and advances the
queue on one button press. A single ATmega32 coordinates both stations, driving an LM35
temperature sensor, an XD-58C optical pulse sensor, a $4\times8$ alphanumeric matrix
keypad and two $16\times2$ character LCDs across 23 of its 32 general-purpose I/O lines.
Unlike the IoT-first architectures that dominate prior work on embedded health monitoring,
the system needs no network, no companion smartphone and no host PC, so a connectivity
outage cannot stall a consultation. Bench validation against commercial reference
instruments places temperature readings within $\pm 1^{\circ}$C and resting pulse within
$\pm 3$~BPM, with button-to-display latency below 1.2~s and queue ordering preserved under
sustained cycling. We also report the clock constraint that shaped the design. At the
factory 1~MHz clock the baud generator cannot reach 38400~bps within the datasheet's error
tolerance, and fuse-selecting the internal 8~MHz oscillator brings the error down to
$0.16\%$.
